\uuid{9rZS}
\titre{Noyau/Image/Injection/Surjection/Bijection/Théorème du rang}
\auteur{hanani ; mzari}
\organisation{exo7}
\licence{CC BY-NC-SA 4.0}
\datecreate{2020-01-20}
\niveau{L1}
\module{Algèbre}
\chapitre{Application linéaire}
\sousChapitre{Définition}
\difficulte{2}
\theme{Noyau et image, Noyau, Image, Injection, Surjection, Bijection, Théorème du rang}
\exo7qcmid{lille-2.applications-lineaires.014}

\contenu{
\begin{question}
\qclassification{107.01}
\qtags{motcle=Noyau/Image/Injection/Surjection/Bijection/Théorème du rang}
Soit $f$ une application linéaire de $\Rr^3$ dans $\Rr^5$. Quelles sont les assertions vraies ?
\begin{answers}
\bad{Si $\ker f = \{(0,0,0)\}$, alors $f$ est surjective}.
\good{Si $\ker f $ est une droite vectorielle, alors $\Im f $ est un plan vectoriel}.
\good{$f$ est injective si seulement si $\dim \Im f =3$}.
\bad{$f$ est  bijective si et seulement si $\ker f=\{(0,0,0)\}$}.
\end{answers}
\begin{explanations} $f$ est injective si et seulement si $\ker f = \{(0,0,0)\}$. $f$ ne peut pas être surjective, puisque d'après le théorème du rang, la dimension de $\Im f $ est au plus $3$.
\end{explanations}
\end{question}
}
